% Figure for Logic and Proofs §11.
\begin{tikzpicture}[x=2.6cm, y=0.62cm,
    dot/.style={circle, fill=black, inner sep=1.1pt},
    num/.style={font=\scriptsize},
    ar/.style={-{Stealth[length=4pt]}, thin}]
  % columns: N_3 (domain of g), N_4 (domain of f), N_5 (codomain of f)
  \foreach \i in {1,2,3} \node[dot] (a\i) at (0,-\i) {};
  \foreach \i in {1,2,3,4} \node[dot] (b\i) at (1,-\i) {};
  \foreach \i in {1,2,3,4,5} \node[dot] (c\i) at (2,-\i) {};
  \foreach \i in {1,2,3} \node[num, left=2pt] at (a\i) {$\i$};
  \foreach \i in {1,3,4} \node[num, above left=0pt and 1pt] at (b\i) {$\i$};
  \node[num, figR, above left=0pt and 1pt] at (b2) {$i_0 = 2$};
  \foreach \i in {1,2,3,4,5} \node[num, right=2pt] at (c\i) {$\i$};
  % g skips i_0 = 2
  \draw[ar, figB] (a1) -- (b1);
  \draw[ar, figB] (a2) -- (b3);
  \draw[ar, figB] (a3) -- (b4);
  % f, with f(i_0) = k+1 = 5
  \draw[ar] (b1) -- (c2);
  \draw[ar, figR, thick] (b2) -- (c5);
  \draw[ar] (b3) -- (c1);
  \draw[ar] (b4) -- (c4);
  % brace marking N_k inside N_{k+1}
  \draw[decorate, decoration={brace, amplitude=4pt}, figG] (2.18,-0.8) -- (2.18,-4.2)
    node[midway, right=5pt, num, figG] {$\mathbb{N}_4$};
  % labels
  \node[font=\small, figB] at (0,-5.6) {$\mathbb{N}_{3}$};
  \node[font=\small] at (1,-5.6) {$\mathbb{N}_{4}$};
  \node[font=\small] at (2,-5.6) {$\mathbb{N}_{5}$};
  \draw[ar, figB] (0.2,-5.6) -- node[above, num] {$g$} (0.8,-5.6);
  \draw[ar] (1.2,-5.6) -- node[above, num] {$f$} (1.8,-5.6);
\end{tikzpicture}
